\begin{ouverture}[espace]{18}{Solides : cubes, pavés et patrons}{Partie 4 · Espace et géométrie}
{Descriptions · Patrons du cube et du pavé · Assemblages}
\objectif{Décrire les solides usuels par leurs faces, arêtes et sommets.}
\objectif{Reconnaître et construire des patrons du cube, reconnaître ceux du pavé.}
\objectif{Dénombrer des cubes et coder des déplacements.}
\prerequis{Solides et patron du cube vus au CM1.}
\end{ouverture}
\pagestyle{otMain}\essentiel\fontsize{11}{15}\selectfont
\phantomsection\label{cours-18-1}
\essSection{1}{Décrire sans se tromper}
\begin{definition}Les solides usuels : cube, pavé, boule, cylindre, cône, pyramide, prisme droit. Un polyèdre n’a que des faces plates ; on le décrit par le nombre et la nature de ses faces, arêtes et sommets.\end{definition}
\label{fig-cours-18-1-cours-solides}\figureGeo{cours-solides}

\begin{propriete}Cube : 6 faces carrées, 12 arêtes, 8 sommets. Pavé : 6 faces rectangulaires, faces opposées identiques. Pyramide : une base et des faces latérales triangulaires. Prisme droit : deux bases superposables reliées par des rectangles.\end{propriete}
\begin{attention}En perspective, les faces cachées existent : les arêtes invisibles se dessinent en pointillés. Justifier un nom de solide par ses propriétés, jamais par l’allure du dessin.\end{attention}
\clearpage
\phantomsection\label{cours-18-2}
\essSection{2}{Patrons du cube et du pavé}
\begin{definition}Un patron, mis à plat d’un seul tenant, se replie en solide sans superposition. Un patron de cube réunit 6 carrés ; un patron de pavé réunit 6 rectangles, identiques deux à deux.\end{definition}
\label{fig-cours-18-2-cours-patron-pave}\figureGeo{cours-patron-pave}

\begin{savoirfaire}[Les trois paires de faces]Tester un patron de pavé : repérer les trois paires de faces (dessus-dessous, devant-derrière, gauche-droite) ; les faces d’une même paire ne se touchent pas sur le patron, et les dimensions des bords pliés ensemble coïncident.\end{savoirfaire}
\begin{exemple}Dans un patron de cube en croix, les faces opposées sont séparées par une face. Sur un patron de pavé, un rectangle de 5 sur 3 ne peut pas se plier contre un bord de 2 : les longueurs des bords communs doivent être égales.\end{exemple}

\clearpage\rubrique{Découverte}
\begin{blocLivre}{ex-18-1}\begin{exercice}[1]{Nommer les solides}
Nommer les six solides dessinés et entourer ceux qui sont des polyèdres.
\label{fig-ex-18-1}\figureGeo{sol-1}
\end{exercice}\end{blocLivre}
\begin{blocLivre}{ex-18-2}\begin{exercice}[1]{Compter sur le cube}
Donner le nombre de faces, d’arêtes et de sommets du cube, puis du pavé.
\label{fig-ex-18-2}\figureGeo{sol-2}
\end{exercice}\end{blocLivre}
\begin{blocLivre}{ex-18-3}\begin{exercice}[1]{Le prisme à base triangulaire}
Décrire le prisme droit dessiné : bases, autres faces, arêtes et sommets.
\label{fig-ex-18-3}\figureGeo{sol-3}
\end{exercice}\end{blocLivre}
\begin{blocLivre}{ex-18-4}\begin{exercice}[1]{La pyramide}
La pyramide dessinée a une base carrée. Donner le nombre de ses faces, arêtes et sommets.
\label{fig-ex-18-4}\figureGeo{sol-4}
\end{exercice}\end{blocLivre}
\begin{blocLivre}{ex-18-5}\begin{exercice}[1]{Patron de cube ?}
Parmi les trois dispositions de six carrés, lesquelles sont des patrons de cube ?
\label{fig-ex-18-5}\figureGeo{sol-5}
\end{exercice}\end{blocLivre}

\rubrique{Entraînement}
\begin{blocLivre}{ex-18-6}\begin{exercice}[2]{Les faces opposées}
Sur le patron en croix, quelle face se retrouve opposée à la face A ? Et à la face E ?
\label{fig-ex-18-6}\figureGeo{sol-6}
\end{exercice}\end{blocLivre}
\begin{blocLivre}{ex-18-7}\begin{exercice}[2]{Patron de pavé ?}
Parmi les deux dispositions de six rectangles, laquelle est un patron de pavé ? Justifier par les paires de faces.
\label{fig-ex-18-7}\figureGeo{sol-7}
\end{exercice}\end{blocLivre}
\begin{blocLivre}{ex-18-8}\begin{exercice}[2]{Le patron du dé}
Les faces opposées d’un dé totalisent 7. Compléter le patron : 1, 2 et 3 sont placés. Où écrire 4, 5 et 6 ?
\label{fig-ex-18-8}\figureGeo{sol-8}
\end{exercice}\end{blocLivre}
\begin{blocLivre}{ex-18-9}\begin{exercice}[2]{Construire le patron}
Tracer un patron d’un pavé de 4 cm sur 2 cm sur 1 cm. Combien de rectangles de chaque taille faut-il ?
\label{fig-ex-18-9}\figureGeo{sol-9}
\end{exercice}\end{blocLivre}
\begin{blocLivre}{ex-18-10}\begin{exercice}[2]{Les cubes cachés}
Compter les cubes de l’assemblage, cubes cachés compris.
\label{fig-ex-18-10}\figureGeo{sol-10}
\end{exercice}\end{blocLivre}

\rubrique{Problèmes}
\begin{blocLivre}{ex-18-11}\begin{exercice}[3]{La bonne vue}
Associer à l’assemblage sa vue de face et sa vue de dessus parmi les trois vignettes.
\label{fig-ex-18-11}\figureGeo{sol-11}
\end{exercice}\end{blocLivre}
\begin{blocLivre}{ex-18-12}\begin{exercice}[3]{Compléter le grand cube}
L’assemblage compte 20 cubes. Combien en manque-t-il pour former un grand cube de 3 cubes de côté ?
\label{fig-ex-18-12}\figureGeo{sol-12}
\end{exercice}\end{blocLivre}
\begin{blocLivre}{ex-18-13}\begin{exercice}[3]{Coder le chemin}
Écrire la suite d’instructions menant le pion du départ au but en suivant le chemin dessiné (avancer, pivoter à droite ou à gauche).
\label{fig-ex-18-13}\figureGeo{sol-13}
\end{exercice}\end{blocLivre}
\begin{blocLivre}{ex-18-14}\begin{exercice}[3]{Exécuter le programme}
Le pion part de la case marquée, tourné vers la droite de la feuille. Exécuter : avancer de 4 ; pivoter à gauche ; avancer de 2 ; pivoter à gauche ; avancer de 4 ; pivoter à gauche ; avancer de 2. Où arrive-t-il ?
\label{fig-ex-18-14}\figureGeo{sol-14}
\end{exercice}\end{blocLivre}
\begin{blocLivre}{ex-18-15}\begin{exercice}[3]{Le cube peint}
Un cube de 3 sur 3 sur 3, fait de 27 petits cubes, est peint sur toutes ses faces. Combien de petits cubes ont exactement 2 faces peintes ? Et aucune ?
\label{fig-ex-18-15}\figureGeo{sol-15}
\end{exercice}\end{blocLivre}

\clearpage
